\chapter{The Normal Distribution}

The Normal distribution, also known as the Gaussian distribution, is a
type of continuous probability distribution for a real-valued random
variable. It is one of the most important probability distributions in
statistics due to its several unique properties and usefulness in many
areas.\index{normal distribution}

\section{Defining the Normal Distribution}

The Normal distribution is defined by its mean ($\mu$) and standard
deviation ($\sigma$). The probability density function (pdf) of a
Normal distribution is given by:

\begin{equation*}
f(x|\mu, \sigma^2) = \frac{1}{\sqrt{2\pi\sigma^2}} \exp\left(-\frac{(x-\mu)^2}{2\sigma^2}\right)
\end{equation*}

where:
\begin{itemize}
\item $x$ is the point up to which the function is integrated,
\item $\mu$ is the mean or expectation of the distribution,
\item $\sigma$ is the standard deviation,
\item $\sigma^2$ is the variance.
\end{itemize}

\section{Importance of the Normal Distribution}
There are several reasons why the Normal distribution is crucial in statistics:

\begin{itemize}
\item \textbf{Central Limit Theorem:} One of the main reasons for the
  importance of the Normal distribution is the Central Limit Theorem
  (CLT). The CLT states that the distribution of the sum (or average)
  of a large number of independent, identically distributed variables
  approaches a Normal distribution, regardless of the shape of the
  original distribution.

\item \textbf{Symmetry:} The Normal distribution is symmetric, which
  simplifies both the theoretical analysis and the interpretation of
  statistical results.

\item \textbf{Characterized by Two Parameters:} The Normal
  distribution is fully characterized by its mean and standard
  deviation. The mean determines the center of the distribution, and
  the standard deviation determines the spread or girth of the
  distribution.

\item \textbf{Common in Nature:} Many natural phenomena follow a
  Normal distribution. This includes characteristics like people's
  heights or IQ scores, measurement errors in experiments, and many
  others.
\end{itemize}

Given its properties, the Normal distribution serves as a foundation
for many statistical procedures and concepts, including hypothesis
testing, confidence intervals, and linear regression analysis.

\subsection{Example}

A set of course grades is modeled by a normal distribution with mean $\mu = 75$
and standard deviation $\sigma = 8$. Using a 90--80--70--60 grading scale,
find the percentage of students who earned a B (scores from 80 to less than 90)
using Excel and Python.

\textbf{Solution}

Let $X$ be a student's grade. Then
\[
X \sim N(75,8).
\]

A B corresponds to the interval $80 \le X < 90$, so we want
\[
P(80 \le X < 90).
\]

Excel computes cumulative probabilities using

\[
\texttt{NORM.DIST(x, mean, standard\_dev, TRUE)}
\]

Therefore,
\[
P(80 \le X < 90)
=
P(X < 90) - P(X < 80).
\]

In Excel, enter:
\[
\texttt{=NORM.DIST(90,75,8,TRUE) - NORM.DIST(80,75,8,TRUE)}
\]

Excel returns approximately
\[
0.2356.
\]

So about $23.56\%$ of students earned a B.

Now compute the same probability in Python.

\begin{minted}{python}
from scipy.stats import norm

mu = 75
sigma = 8

probability = norm.cdf(90, mu, sigma) - norm.cdf(80, mu, sigma)

print(probability)
print(probability * 100)
\end{minted}

Python returns approximately

\[
0.2356
\]

which confirms the Excel result.

Both Excel and Python are computing the same cumulative distribution function.

\begin{Exercise}[label = normalt1]

A batch of parts has diameters modeled by a normal distribution with mean
$\mu = 10$ mm and standard deviation $\sigma = 2$ mm:
\[
X \sim N(10,2).
\]
A part is considered \emph{acceptable} if its diameter is between $6.72$ mm and
$12.28$ mm.

\begin{itemize}
    \item[(a)] Write the probability statement that represents the fraction of acceptable parts.
    \item[(b)] Use Excel to compute this probability using \texttt{NORM.DIST}.
    \item[(c)] Convert your probability to a percentage.
\end{itemize}

\vspace{100mm}

\end{Exercise}

\begin{Answer}[ref = normalt1]

\textbf{(a)} The fraction of acceptable parts is
\[
P(6.72 \le X \le 12.28).
\]

\textbf{(b)} Since \texttt{NORM.DIST(x, mean, sd, TRUE)} returns $P(X \le x)$,
we compute
\[
P(6.72 \le X \le 12.28)=P(X \le 12.28)-P(X \le 6.72).
\]

In Excel, enter:
\[
\texttt{=NORM.DIST(12.28,10,2,TRUE) - NORM.DIST(6.72,10,2,TRUE)}.
\]

\textbf{(c)} Excel returns approximately $0.8224$, so the percentage acceptable is
\[
82.24\%.
\]

\end{Answer}
